
In order to validate of the proposed design, we simulated a model of our robot with V-REP simulation framework~\cite{vrep13}. Here, we kept the size of the vehicle and estimated the real mass and inertia of the fully loaded robot. The design parameters used are the ones shown in Table~\ref{tab:selected}.  The vehicle mass is 6.05kg. In V-REP we simulated the force and torque produced by each propeller as modeled in~\eqref{eq:Fi} and~\eqref{eq:Mi}. Thus, this simulation includes the validation of both the design discussed in Section~\ref{sec:design-parameters} and the controller proposed in Section~\ref{sec:posit-attit-contr}.

The validation was divided in two phases: on the first, each motion mode (translation and attitude) is evaluated separately, and on the second, we simulated a path defined by a sequence of waypoints with attitude setpoints and a unmodeled payload transportation, thus combining both modes simultaneously. In the last case we added noise to the pose with a variance of 2cm on position, 2cm/s on velocity, 5\degree \ on attitude and 1\degree/s on angular velocity. 

\subsection{Individual motion mode}

% The analysis of each mode enables a better adjustment of the gains of the controller and validates each controller, and component, individually. 
For this validation, the robot is holding on position $(0,0,1)$, while the reference is set at $(1,2,4)$. The total movement of the robot is 1, 2 and 3 meters along X, Y and Z axis. On Figure~\ref{fig:position_errors} we provide the obtained results for position convergence with the proposed controller in~\eqref{eq:fctrl} and~\eqref{eq:pctrl}. Convergence rates are equal across all axis as expected. It only depends on the gains (equal for all axis) and the mass of the vehicle.

\begin{figure}
  \centering
  \includegraphics[width=0.8\linewidth]{error-p.pdf}
  \caption{Translational mode error convergence, from top to bottom: $e_x$ and $e_v$.}
  \label{fig:position_errors}
\end{figure}

On the attitude simulation, we start the vehicle at $(0\degree,0\degree,0\degree)$ and the reference is set as $(70\degree,-50\degree,30\degree)$ (both in $XYZ$ Euler angles). The results are shown in Figure~\ref{fig:attitude_errors}. Attitude convergence is approximately the same on all axis. As it depends only on the gains (equal for all axis) and the inertia of the vehicle (almost diagonal due to its symmetry), this result was also expectable.

\begin{figure}
  \centering
  \includegraphics[width=0.8\linewidth]{error-a.pdf}
  \caption{Attitude mode error convergence, from top to bottom: $e_R$ and $e_\omega$.}
  \label{fig:attitude_errors}
\end{figure}


\subsection{Waypoint navigation and payload transportation}

To validate the joint position and attitude control, we used waypoint navigation: we created a virtual path composed by 6 waypoints, varying on both position and attitude. The trajectory followed by the vehicle is represented on \ref{fig:trajectory}. Plotted are both the X and Y axes of the body frame~$\mathcal{B}$ to represent its attitude. We retrieved the errors across the whole trajectory following, for both position and atitude. These are shown in Figure~\ref{fig:nav_errors}.

\begin{figure}
  \centering
  \includegraphics[width=0.75\linewidth]{3d-trajectory.pdf}
  \caption{Trajectory followed by the robot without load. Plotted are X and Y axis of the inertia frame~$\mathcal{I}$.}
  \label{fig:trajectory}
\end{figure}

\begin{figure}
  \centering
  \includegraphics[width=0.8\linewidth]{robust-noload.pdf}
  \caption{Position (on top) and attitude (on bottom) errors along the no-load trajectory of Figure \ref{fig:trajectory}.}
  \label{fig:nav_errors}
\end{figure}

Then, we added a non-modeled payload to the system: a sphere with 6Kg of mass --- about the same mass of the vehicle. Even with such payload, the controller was able to converge to the required positions, though taking significantly more time. Shown on Figure~\ref{fig:loaded} are screenshots of the vehicle along its 3 waypoint trajectory: 1~meter up along Z and 1~meter right along Y. The position and attitude errors for this trajectory are plotted on Figure \ref{fig:nav_errors_load}

\begin{figure}
  \centering
  \includegraphics[width=0.75\linewidth]{snapshots.jpg}
  \caption{Trajectory followed by the robot with 6kg non-modeled load. Pictured are six positions of the robot along its trajectory.}
  \label{fig:loaded}
\end{figure}

\begin{figure}
  \centering
  \includegraphics[width=0.8\linewidth]{robust-load.pdf}
  \caption{Position (on top) and attitude (on bottom) errors with non-modeled 6kg load along the trajectory of Figure \ref{fig:loaded}.}
  \label{fig:nav_errors_load}
\end{figure}

%%% Local Variables:
%%% mode: latex
%%% TeX-master: "main"
%%% End:
